% !TEX root = ../main.tex
% ============================================================================
% Chapter 6: 总结与展望
% 作用: 全文综合（不再引入新实验），归纳结论 / 局限 / 未来工作
% ============================================================================

\section{总结与展望}

\subsection{工作总结}

本文围绕NSL-KDD数据集展开网络入侵检测的系统研究，主要工作可归纳为以下五个方面。

第一，深入分析了NSL-KDD数据集的结构特征与建模挑战。该数据集训练集包含125,973条样本，测试集22,544条样本，标签涵盖Normal与四种大类攻击（DoS、Probe、R2L、U2R）。测试集中存在14种训练阶段从未出现的新型攻击（如apache2、sqlattack、worm等），这一设计虽然有助于评估模型的开放世界泛化能力，但也成为多分类任务准确率难以突破的上限约束。

第二，构建了完整的特征工程流水线。通过OneHot编码将41维原始特征扩展为53维，经方差阈值过滤降至32维，再借助随机森林重要度筛选压缩至20维关键特征。Top-5特征包括flag\_SF、same\_srv\_rate、dst\_host\_srv\_count、dst\_host\_same\_srv\_rate、diff\_srv\_rate等连接状态与流量统计指标，证明树模型筛选的特征在DT与RF之间具有跨模型族的稳定性排序。

第三，系统评估了传统机器学习方法。决策树在统一测试集上取得最高F1值0.7621与最快推理速度（0.02秒/批次），随机森林则以0.9150的AUC在概率排序能力上领先，两者形成"高精度硬决策"与"高判别软输出"的互补格局。

第四，对比了三种神经网络架构。MLP在IID验证集上F1达到0.989，但在含14种unseen攻击的OOD测试集上骤降至0.7202，验证测试差距高达26个百分点。CNN与LSTM的表现与MLP接近但未实现反超，印证了"模型越复杂效果越好"的朴素直觉在表格型网络流量数据上并不成立。

第五，多分类任务揭示了类别不平衡与小样本问题的双重制约。MLP的full\_accuracy为0.1009，约为RF与DT的两倍，但三者的f1\_macro均低于0.02。SMOTE过采样实验呈现显著负向效果，多分类准确率从10.09\%降至2.84\%，降幅高达72\%，说明在小类（U2R仅52条样本）上合成新样本会污染训练分布，反而加剧误分类。

综合而言，DT与RF在二分类任务中仍是工程实践的最优选择；深度学习方法的IID优势无法迁移至OOD场景；多分类任务需要跳出传统监督分类的框架才能进一步突破。

\subsection{研究局限}

本研究的局限性主要体现在以下六个方面。

其一，多分类任务受限于训练与测试之间的类别分布错位。训练集仅含23种具体攻击类型，测试集则含38种，测试集中15种未见攻击从根本上限制了全量准确率上限，所有模型在该任务上的绝对数值偏低，并非模型缺陷，而是数据集的固有特性。

其二，验证集指标显著高估了真实泛化能力。MLP在IID验证集上的F1为0.989，但同一模型在OOD测试集上的F1仅为0.7202。这一现象提示研究者不应仅依赖验证集表现评估模型，分布外泛化能力才是入侵检测真实部署场景下的关键指标。

其三，实验未进行统计显著性检验。6个模型的二分类F1差异较小（DT 0.7621 vs MLP 0.7202），但未通过McNemar检验或配对t检验验证差异是否显著，结论的统计严谨性有待加强。

其四，CNN与LSTM的架构选择与数据特性不匹配。NSL-KDD每条样本代表一次独立的网络连接记录，样本之间无空间局部结构或时序依赖关系，CNN的平移不变性与LSTM的时序建模能力均缺乏发挥空间，这可能是两者未能显著超越MLP的深层原因。

其五，未系统测量推理延迟与资源消耗。安全运营场景对实时性要求较高，但本研究仅在测试集上做了批量推理，未评估单样本推理延迟、内存占用、CPU/GPU利用率等部署相关指标，难以直接支撑上线决策。

其六，实验仅基于NSL-KDD单一数据集。NSL-KDD作为1999年KDD Cup的改进版，距今已二十余年，流量模式与攻击类型与现代网络环境存在显著差距，研究结论在CICIDS2017、UNSW-NB15等现代数据集上的可迁移性尚未验证。

\subsection{未来工作}

基于上述结论与局限，未来研究可从以下六个方向展开。

第一，引入开放集识别（Open-set Recognition）框架处理unseen攻击。通过在传统分类器输出端增加开放集判别模块（如OpenMax、Energy-based Score），将已知攻击与新型攻击区分开来，使多分类任务摆脱封闭集假设的限制。

第二，使用Focal Loss或class\_weight参数替代SMOTE缓解类别不平衡。Focal Loss通过降低易分样本的损失权重，引导模型关注U2R、R2L等小类；class\_weight则在损失函数层面直接放大少数类梯度。二者均无需合成新样本，可避免SMOTE实验中观察到的训练分布污染问题。

第三，探索集成学习方法（Stacking、Voting）融合各模型优势。DT的高精度与快推理、RF的优概率排序、MLP的高已知类识别准确率具有互补性，通过元学习器或加权投票集成，有望在单一指标上超越任一基模型。

第四，在CICIDS2017、CSE-CIC-IDS2018、UNSW-NB15等现代数据集上验证结论的普适性。这些数据集包含Web攻击、慢速DDoS等新型威胁，更贴近真实网络环境，跨数据集验证可揭示当前结论的适用范围与边界条件。

第五，面向实时入侵检测部署需求开展模型压缩与推理优化。通过量化（Quantization）、剪枝（Pruning）、知识蒸馏（Knowledge Distillation）等技术压缩DT/RF/MLP体积，在边缘设备或高吞吐链路上实现毫秒级单样本推理，并系统测量延迟、吞吐量与资源消耗。

第六，探索图神经网络（GNN）建模网络流量拓扑。将网络节点、连接关系作为图结构输入，借助GNN的消息传递机制捕捉节点间的拓扑特征，可能突破传统表格模型对独立同分布假设的依赖，为入侵检测提供新的归纳偏置。

最后，在线学习与增量更新机制也是值得探索的方向。真实网络环境中攻击模式持续演化，模型需要具备从新样本中持续学习的能力，同时避免灾难性遗忘。这一方向与开放集识别、主动学习共同构成下一代入侵检测系统的核心研究议题。
